% Page 048 — page imprimée 43
% Contenu seul : le préambule, l'en-tête et la pagination sont gérés par meyrueis.tex

Quant au pasteur Guiraud, incomparable orateur, doué d’une voix plus distinguée que
forte, quel enchantement que d’entendre cet homme ! Je me souviens d’une prière de
lui, entendue à Meyrueis, qui commençait par cette invocation : « \emph{O Toi qui donnes,
Dieu de Noël, O Toi qui pardonnes, Dieu de Pâques, nous venons à Toi !} [\ldots]

Pendant deux ans, les circonstances locales s’y prêtant, je fis la classe de 6\textsuperscript{e} et de 5\textsuperscript{e} à
une jeune fille très douée, dont les parents, excellents paroissiens, ne voulaient pas se
séparer. Elle s’appelait Jeanine Fabre.
Je me fis donc professeur. Ma femme et moi prîmes en charge toutes les matières du
programme : français, latin, anglais, histoire, géographie, botanique géologie etc. La
suite de ces leçons elle fut admise avec félicitations dans un lycée de la capitale. Outre
la joie que l’on éprouve quand on a à faire à une enfant intelligente, cela apporta un
peu d’aisance à notre foyer qui en avait bien besoin. La contre-partie, un peu fâcheuse
c’est que toutes nos matinées furent prises. Autre conséquence : un bruit favorable se
répandit autour de ces leçons, et des vacanciers voulurent également que, au cours de
l’été, je fasse travailler leurs enfants. Je ne m’y suis pas refusé.

Janine Fabre, après y avoir fait ses études, fut nommée professeur de lettres au Lycée
Victor Duruy dans le 7\textsuperscript{e} Arrondissement.

\begin{itemize}
\item \emph{Plus tard elle épousa Francis Nougarède, ingénieur géologue. Ils eurent trois enfants,
Florence, Olivier et Pascal qui revenaient ensemble ou à tour de rôle au « Clair Logis »
leur maison de Meyrueis qui avait auparavant appartenu à la tante de Janine, le Dr. Juliette
Cabanel,
Francis fut emporté par une tumeur en 2004 et Janine le suivit un an plus tard.}
\end{itemize}

\emph{« Le Clair Logis » est passé désormais en d’autres mains.}

Nous avions aussi au mois de juillet et d’août étudiants en Théologie ou en Lettres
comme pensionnaires. Au nombre de 3 ou 4 chaque fois. Ils nous étaient envoyés par
l’Association Libérale de Strasbourg. Nous les faisions coucher dans deux vastes
chambres au premier. Toute la journée ils couraient dans les montagnes, visitaient la
région et, le soir venu, ils entonnaient des cantiques au presbytère. Cela tenait lieu de
culte de famille. Quel ne fut pas notre étonnement de constater que des
rassemblements assez importants de touristes se formaient autour de notre demeure
pour les écouter chaque jour ! Ces voix d’hommes étaient très belles et le sens musical
de ces jeunes strasbourgeois incontestable !

Comment ne pas mentionner nos rapports avec le Docteur Lapeyre et mes escapades
en sa compagnie sur les Causses. Dans une petite cité les rapports entre le pasteur et
Le médecin sont très importants, à tous égards. A Meyrueis, le docteur, catholique
pratiquant, préférait la compagnie du pasteur à celle du prêtre. Il l’avait déjà montré
avec un de mes prédécesseurs, M. Maurin.
Sa clientèle, en dehors du centre et des vallées, s’étendait sur les Causses, territoire
dont les habitants m’étaient inconnus, car je n’y avais aucun paroissien. Son rayon
d’action atteignait la Malène et Sainte Enimie, sur le Tarn, juste avant les Gorges
proprement dites. n’aimait pas rouler seul dans sa voiture, une belle conduite
intérieure Citroën confortable.
